\section{Results}\label{sec:result_analysis}

\subsection{RQ1: MR-based Behavioral Coverage}\label{sec:results:rq1}
We first examine the gap between behavioral reach and behavioral validation in existing component test suites.
Table~\ref{tab:completeness} reports MR completeness across libraries and component categories using the \textit{Touch} and \textit{Cover} metrics.

\subsubsection{Existing tests achieve high behavioral reach but substantially lower behavioral validation.} 
Table~\ref{tab:completeness} shows that \textit{Touch} is consistently higher than \textit{Cover} across libraries, categories, and LLM configurations\footnote{DeepSeek (\textit{deepseek-chat}), Gemini (\textit{gemini-2.5-flash}), and GPT (\textit{gpt-5-mini}).}. At the model level, the overall MR \textit{Cover} rate ranges from 42.5\% to 47.6\% across the three LLM configurations, while MR \textit{Touch} remains substantially higher.
Using \textit{Gemini}, library-level \textit{Touch} ranges from 85.4\% to 89.4\%, whereas \textit{Cover} ranges from 41.9\% to 54.7\%. This gap suggests that existing tests often exercise the inputs, states, or interactions associated with inferred MRs, but do not always include assertions that validate the expected behavioral relation. Thus, within the inferred MR space, many uncovered relations arise not because the associated behavior is absent from tests, but because the tests do not provide explicit relation-level evidence for the expected behavior.

\begin{table}[h]
\centering
\caption{MR completeness by UI library and component category.}
\label{tab:completeness}
\setlength{\tabcolsep}{2pt}
% \renewcommand{\arraystretch}{1.1}
\begin{tabular}{l l ccc ccc}
\toprule
 &  & \multicolumn{3}{c}{\textbf{Touch} (\%)} & \multicolumn{3}{c}{\textbf{Cover} (\%)} \\
\cmidrule(lr){3-5} \cmidrule(lr){6-8}
\textbf{Group} & \textbf{Item} & Gemini & DeepSeek & GPT & Gemini & DeepSeek & GPT \\
\midrule
\multirow{4}{*}{Library} 
 &  ant-design                           & 88.2 & 83.0 & 82.7 & 45.0 & 40.9 & 38.1 \\
 & element-plus                        & 85.4 & 78.3 & 81.2 & 54.7 & 50.3 & 50.5 \\
 &  base-ui                             & 89.4 & 84.4 & 86.8 & 45.7 & 43.7 & 40.8 \\
 & material-ui                         & 88.6 & 78.0 & 83.1 & 41.9 & 41.5 & 37.6 \\
\midrule
\multirow{5}{*}{Category}
 &  Data display                        & 84.8 & 74.5 & 79.4 & 40.1 & 36.5 & 36.1 \\
 & Feedback                            & 88.1 & 80.4 & 80.6 & 45.6 & 44.9 & 36.4 \\
 &  Inputs                               & 89.3 & 87.9 & 88.9 & 57.8 & 56.3 & 55.9 \\
 & Layout                              & 76.0 & 60.2 & 67.2 & 37.6 & 22.7 & 25.2 \\
 &  Navigation                           & 93.7 & 86.8 & 87.3 & 47.1 & 46.0 & 41.5 \\
\bottomrule
\end{tabular}
\end{table}


\subsubsection{Completeness varies across libraries and component categories.} At the library level, \textit{element-plus} shows the highest \textit{Cover}, exceeding 50\% across all three models, whereas \textit{material-ui} remains between 37.6\% and 41.9\%. These differences may reflect variation in library-specific testing practices, component designs, or assertion styles.
Variation is also visible across component categories. \textit{Input} components achieve the highest \textit{Cover} of 57.8\%, whereas \textit{Layout} components achieve the lowest (22.7--37.6\%). This pattern is consistent with the nature of the underlying behaviors: input-related relations often expose explicit state transitions, validation outcomes, and callback effects that are straightforward to assert, while layout relations more often involve visual, spatial, or responsive behavior that may be harder to validate through conventional component tests.



\begin{tcolorbox}[colback=gray!5, colframe=gray!30, title=\textbf{\small Answer to RQ1}, coltitle=black!90, left=2.5pt, top=3pt, bottom=3pt, right=2.5pt, boxsep=1.8pt]
Existing UI component test suites achieve high behavioral reach but lower behavioral validation. Across libraries and models, inferred behavioral relations are more often exercised than covered. Validation strength also varies by component category, with \textit{Input} components consistently better covered than \textit{Layout} components.
\end{tcolorbox}


\subsection{RQ2: Behavioral Gaps}\label{sec:rq2-results}

RQ2 examines where behavioral gaps arise among inferred relations that are not covered by existing tests. We first decompose uncovered relations into \textit{Weak} and \textit{Untouched} relations, and then identify the relation types that remain consistently under-covered across LLMs.

\subsubsection{Most uncovered relations are associated with weak validation rather than missing test scenarios.}

Figure~\ref{fig:rq2_gap_decomp} divides inferred relations into \textit{Cover}, \textit{Weak}, and \textit{Untouched}. Across the three LLMs, 52.4--57.5\% of inferred relations remain uncovered under the \textit{Cover} criterion. Weak-oracle cases form the majority of these uncovered relations, accounting for roughly two-thirds to three-quarters of them. In contrast, only 12.5--19.4\% of all inferred relations remain completely untouched. This indicates that most uncovered relations are exercised by existing tests but are not explicitly validated. Thus, the observed gaps are more often associated with limited behavioral validation than with entirely missing test scenarios.


\begin{figure}[h]
  \centering
  \includegraphics[width=0.86\linewidth]{figure/04-rq2-distribution.png}
  \caption{Distribution of MR coverage outcomes (Cover, Weak, Untouched) for three LLMs across 214 components.}
  \label{fig:rq2_gap_decomp}
\end{figure}

\begin{figure}[h]
  \centering
  \includegraphics[width=0.96\linewidth]{figure/05-rq2-relation.png}
  \caption{Relation types with the lowest \textit{Cover} rates, grouped by top-level categories. Colored brackets distinguish the groups (blue: Visual/layout, orange: Input/prop, green: Composition/\
  context, brown: Interaction/accessibility).}
  \label{fig:rq2_worst_types}
\end{figure}

\subsubsection{Visual/layout behavior is prominent among the lowest-covered relation types.} 

Figure~\ref{fig:rq2_worst_types} shows the eight relation types with the lowest \textit{Cover} rates. Visual/layout behavior is the most represented top-level category in this group, covering four of the eight relation types: \textit{overflow handling} (17.4\%), \textit{animation consistency} (18.8\%), \textit{visual alignment} (27.9\%), and \textit{responsive sizing} (34.8\%).
The breakdown further shows that low \textit{Cover} can arise from different sources. For \textit{overflow handling}, the gap is mainly due to missing behavioral reach: 45.0\% of inferred relations are \textit{Untouched}, suggesting that overflow, clipping, and boundary-layout behaviors are often absent from existing test scenarios. In contrast, \textit{animation consistency} and \textit{monotonicity} have high \textit{Weak} rates, at 56.3\% and 53.5\%, respectively. These relations are often exercised but not explicitly validated, indicating a weak-oracle gap rather than a complete absence of test scenarios.

Similar weak-oracle patterns also appear in \textit{context propagation}, \textit{keyboard interaction}, and \textit{focus management}. These behaviors often require assertions over layout changes, visual state, accessibility attributes, focus state, or cross-component interactions, which are less directly captured by conventional rendering or state-transition checks.


\begin{tcolorbox}[colback=gray!5, colframe=gray!30, title=\textbf{\small Answer to RQ2}, coltitle=black!90, left=2.5pt, top=3pt, bottom=3pt, right=2.5pt, boxsep=1.8pt]
MR gaps are dominated by weak validation rather than entirely missing behavioral scenarios. Across LLMs, weak-oracle relations form the majority of uncovered relations, suggesting that existing tests often exercise relevant behaviors without explicitly checking the expected relations. Low-coverage relation types are frequently observed in visual/layout, accessibility, and contextual behaviors.
\end{tcolorbox}


\subsection{RQ3: Complementary Adequacy Signals}\label{sec:rq3-results}
RQ3 investigates whether MR coverage provides a relation-level adequacy signal beyond statement and branch coverage. We analyze this relationship using the 201 components for which runtime coverage data are available.

\subsubsection{MR Coverage Changes with Sampled Test-suite Size.}
To examine how MR-based metrics change as test suites grow, we construct sampled sub-suites containing 25\%, 50\%, and 75\% of the original tests, repeat the sampling 10 times for each size, and compare them with the full test suites. Table~\ref{tab:rq3-merged} shows that all metrics increase with sampled test-suite size. Statement coverage rises from 54.3\% to 92.7\%, branch coverage from 61.0\% to 83.5\%, MR \textit{Touch} from 44.8\% to 82.2\%, and MR \textit{Cover} from 18.8\% to 45.2\%. Thus, MR coverage grows with test-suite size rather than being detached from test execution.

Statement coverage is moderately correlated with MR \textit{Touch} for smaller sampled suites (e.g., $\rho=0.41$ at 25\%), which is expected because \textit{Touch} only requires exercising the corresponding behavior. In contrast, correlations involving MR \textit{Cover} remain weak and approach zero for the full test suites. Thus, although both metric families increase with test-suite size, structural coverage provides only a partial view of behavioral adequacy: MR \textit{Cover} reflects whether exercised behaviors are explicitly validated, not merely executed.


\begin{table}[h]
\centering
\caption{Coverage growth and within-size Spearman correlations under sampled test-suite sizes ($n=201$ components; 10 random samples per size).}
\label{tab:rq3-merged}
\small
\setlength{\tabcolsep}{4pt}
% \renewcommand{\arraystretch}{1.08}
\begin{tabular}{lrrrr}
\toprule
\textbf{Metric} & \textbf{25\%} & \textbf{50\%} & \textbf{75\%} & \textbf{100\%} \\
\midrule
\multicolumn{5}{l}{\textbf{Mean coverage metrics}} \\
Statement coverage (\%) & 54.3 & 68.9 & 78.7 & 92.7 \\
Branch coverage (\%)    & 61.0 & 69.7 & 75.8 & 83.5 \\
MR Touch (\%)           & 44.8 & 58.6 & 66.6 & 82.2 \\
MR Cover (\%)           & 18.8 & 28.1 & 34.5 & 45.2 \\
\midrule
\multicolumn{5}{l}{\textbf{Within-size Spearman correlations}} \\
$\rho$(Statement, MR Cover) & 0.24 & 0.19 & 0.12 & $-$0.16 \\
$\rho$(Branch, MR Cover)    & 0.00 & 0.03 & 0.01 & $-$0.07 \\
$\rho$(Statement, MR Touch) & 0.41 & 0.30 & 0.19 & $-$0.12 \\
$\rho$(Branch, MR Touch)    & 0.07 & 0.08 & 0.05 & 0.02 \\
\bottomrule
\end{tabular}
\end{table}


\subsubsection{Execution Coverage Provides Limited Evidence of Behavioral Validation.} 
We further examine whether the conditional-probability results are sensitive to the \textit{MR Cover} threshold by repeating the analysis with thresholds of 40\%, 50\%, and 60\%. As shown in Table~\ref{tab:rq3-conditional}, the same asymmetry is observed across thresholds. Among components with high statement coverage, the proportion that also reaches the \textit{MR Cover} threshold decreases from 58.2\% to 28.8\% as the threshold becomes stricter. A similar pattern appears for high branch coverage, where the proportion decreases from 61.0\% to 30.1\%. In contrast, components that reach the \textit{MR Cover} threshold more often also have high statement or branch coverage. These results suggest that high \textit{MR Cover} often co-occurs with high execution coverage, but high execution coverage alone does not necessarily indicate high behavioral validation.

\begin{table}[h]
\centering
\caption{Sensitivity of conditional probabilities under different \textit{MR Cover} thresholds. $\mathrm{MRC}_{\theta}$ denotes MR \textit{Cover} $\geq \theta$. Percentages are shown with counts in parentheses.}
\label{tab:rq3-conditional}
\setlength{\tabcolsep}{1.5pt}
\renewcommand{\arraystretch}{1.3}
\begin{tabular}{lrrr}
\toprule
\textbf{Condition} & \textbf{$\theta=40\%$} & \textbf{$\theta=50\%$} & \textbf{$\theta=60\%$} \\
% \textbf{Condition} & \textbf{$\mathrm{MR}_{C\theta}\geq$ 40\%}  & \textbf{$\mathrm{MR}_{C\theta}\geq$ 50\%} & \textbf{$\mathrm{MR}_{C\theta}\geq$ 60\%} \\
\midrule
$P(\mathrm{MRC}_{\theta} \mid \mathrm{Stmt}\geq90\%)$
& 58.2 {\footnotesize (85/146)}
& 41.8 {\footnotesize (61/146)}
& 28.8 {\footnotesize (42/146)} \\

$P(\mathrm{Stmt}\geq90\% \mid \mathrm{MRC}_{\theta})$
& 73.3 {\footnotesize (85/116)}
& 70.1 {\footnotesize (61/87)}
& 65.6 {\footnotesize (42/64)} \\

$P(\mathrm{MRC}_{\theta} \mid \mathrm{Branch}\geq80\%)$
& 61.0 {\footnotesize (83/136)}
& 43.4 {\footnotesize (59/136)}
& 30.1 {\footnotesize (41/136)} \\

$P(\mathrm{Branch}\geq80\% \mid \mathrm{MRC}_{\theta})$
& 71.6 {\footnotesize (83/116)}
& 67.8 {\footnotesize (59/87)}
& 64.1 {\footnotesize (41/64)} \\
\bottomrule
\end{tabular}
\end{table}


\begin{tcolorbox}[colback=gray!5, colframe=gray!30, title=\textbf{\small Answer to RQ3}, coltitle=black!90, left=2.5pt, top=3pt, bottom=3pt, right=2.5pt, boxsep=1.8pt]
MR coverage provides adequacy information beyond execution-based coverage. Although MR metrics increase with sampled test-suite size, weak within-size correlations show that high statement or branch coverage does not necessarily imply high behavioral validation. \textit{MR Cover} specifically captures whether exercised behaviors are explicitly checked against inferred behavioral relations.
\end{tcolorbox}



\subsection{RQ4: Practical Relevance}\label{sec:rq4-results}

RQ4 examines the practical relevance of the inferred MR space and MR-based coverage labels. We evaluate this through three complementary analyses: issue-description mapping, targeted assertion augmentation for weak-oracle relations, and MR-relevant injected fault detection. These analyses are intended to provide supporting evidence rather than causal evidence of real-world fault proneness.

\subsubsection{Reported issue annotations are substantially represented in the inferred MR space.}

% \begin{table}[t]
% \centering
% \caption{Reported issue descriptions mapped to inferred MR relation types.}
% \label{tab:issue-mr-map}
% \setlength{\tabcolsep}{3pt}
% \renewcommand{\arraystretch}{1.08}
% \scriptsize
% \resizebox{\columnwidth}{!}{%
% \begin{tabular}{lrrrr}
% \toprule
% \textbf{MR relation type} & \textbf{Annotated} & \textbf{Mapped} & \textbf{Cov$<$30\%} & \textbf{Cov$\geq$50\%} \\
% \midrule
% Env. adaptation        & 43 & 9/43   & 5/9   & 4/9   \\
% State--visual mapping  & 35 & 23/35  & 10/23 & 13/23 \\
% State synchronization  & 33 & 23/33  & 8/23  & 15/23 \\
% Placement consistency  & 28 & 21/28  & 4/21  & 11/21 \\
% Visual alignment       & 26 & 15/26  & 12/15 & 3/15  \\
% Overflow handling      & 25 & 12/25  & 9/12  & 3/12  \\
% Slot composition       & 22 & 21/22  & 2/21  & 15/21 \\
% Responsive sizing      & 19 & 15/19  & 3/15  & 9/15  \\
% Focus management       & 18 & 16/18  & 9/16  & 6/16  \\
% Event propagation      & 17 & 11/17  & 4/11  & 7/11  \\
% \midrule
% Other mapped types     & 193 & 131/193 & 56/131 & 67/131 \\
% \midrule
% \textbf{All mapped}    & \textbf{459/634} & \textbf{297/459} & \textbf{122/297} & \textbf{153/297} \\
% \bottomrule
% \end{tabular}%
% }
% \end{table}

\begin{table}[h]
\centering
\caption{Summary of reported issue annotations mapped to the inferred MR space.}
\label{tab:issue-mr-summary}
\small
\setlength{\tabcolsep}{4pt}
% \renewcommand{\arraystretch}{1.08}
\begin{tabular}{lrr}
\toprule
\textbf{Outcome} & \textbf{Count} & \textbf{\%} \\
\midrule
Issue annotations & 634 & 100.0 \\
Mapped to MR relation types & 459 & 72.4 \\
\quad Direct mapping & 90 & 14.2 \\
\quad Inferred mapping & 262 & 41.3 \\
\quad Multiple MR mapping & 107 & 16.9 \\
No clear MR match & 175 & 27.6 \\
\bottomrule
\end{tabular}
\end{table}

We first examine whether reported issue descriptions can be related to the inferred MR space. We collected component-related issue reports and manually annotated 634 behavioral descriptions, referred to as issue annotations. Table~\ref{tab:issue-mr-summary} shows that, under a strict criterion requiring a unique MR match, 352 of 634 issue annotations (55.5\%) align with inferred MR relation types. Including annotations with multiple MR matches increases the mapped set to 459 annotations (72.4\%), while the remaining 175 annotations (27.6\%) have no clear match in the current MR space.

Figure~\ref{fig:rq4_issue_scatter} further relates mapped issue annotations to average MR \textit{Cover} across relation types. The mapped issues are distributed across both higher and lower-cover relation types. Frequently mapped relations such as state synchronization and state--visual mapping show relatively high \textit{Cover}, whereas visual alignment, overflow handling, theme consistency, and keyboard interaction fall below the 40\% visual reference line. These results suggest that the inferred MR space represents many behavioral concerns described in reported issues, and that MR \textit{Cover} helps indicate which issue-aligned areas remain less explicitly validated.

\begin{figure}[h]
  \centering
  \includegraphics[width=0.84\linewidth]{figure/07-issue_scatter.png}
  \caption{Relationship between mapped issue annotations and average MR Cover across MR relation types. The dashed line at 40\% is a visual reference for lower-cover relation types, not a decision threshold.}
  \label{fig:rq4_issue_scatter}
\end{figure}


\subsubsection{Weak-oracle relations highlight opportunities for additional behavioral validation.}

As a second perspective on practical relevance, we conducted the oracle-strengthening study described in Section~\ref{sec:external-validation} to examine whether weak-oracle relations correspond to opportunities for strengthening existing tests through targeted assertions derived from the inferred MRs.


\begin{table}[h]
\centering
\caption{Weak-oracle assertion study summary. Percentages are computed with respect to the denominator shown in each block.}
\label{tab:wo-study}
\setlength{\tabcolsep}{3pt}
% \renewcommand{\arraystretch}{1.08}
\begin{tabular}{@{}lcc@{}}
\toprule
\textbf{Outcome} & \textbf{Count} & \textbf{\%} \\
\midrule
\multicolumn{3}{@{}l}{\textbf{Sampled weak-oracle relations} ($n=200$)} \\
\quad Executed with augmented assertion & 200 & 100.0 \\
\midrule
\multicolumn{3}{@{}l}{\textbf{Assertion result} ($n=200$)} \\
\quad Passed -- assertion satisfied & 81 & 40.5 \\
\quad Failed -- assertion not satisfied & 119 & 59.5 \\
\midrule
\multicolumn{3}{@{}l}{\textbf{Failure inspection} ($n=119$)} \\
\quad Mismatch with inferred MR assertion & 84 & 70.6 \\
\quad Snapshot artifact & 14 & 11.8 \\
\quad Indeterminate & 21 & 17.6 \\
\midrule
\textbf{Conclusive outcomes} ($81+84$ of 200) & \textbf{165} & \textbf{82.5} \\
\bottomrule
\end{tabular}
\end{table}


Table~\ref{tab:wo-study} summarizes the results. All 200 augmented tests executed successfully. Among them, 81 added assertions passed, whereas 119 failed. Manual inspection attributed 84 of the failed cases to mismatches between the observed behavior and the inferred MR assertion; the remaining failures resulted from snapshot artifacts or indeterminate execution conditions. These mismatches should not be interpreted as confirmed defects or as direct evidence that the inferred MR is correct. Rather, they identify cases where the existing tests do not provide sufficient oracle evidence to determine whether the inferred relation holds.
Overall, 165 cases (82.5\%) produced informative outcomes, either because the added assertion was satisfied or because manual inspection identified an MR-related mismatch. These results suggest that many weak-oracle relations correspond to behaviors that are exercised but not explicitly validated by existing tests. More broadly, targeted assertion augmentation provides additional behavioral evidence for identifying opportunities to strengthen behavioral validation in existing tests.


\subsubsection{MR labels show a trend in behavior-oriented fault detection.}
To further examine this trend in our targeted injected-fault sample, we conducted the injected-fault study described in Section~\ref{sec:external-validation}. This study is not intended to estimate general fault-detection capability; rather, it provides a focused check of whether MR labels are consistent with detection outcomes for faults explicitly constructed to violate the corresponding MR.
Table~\ref{tab:mr_fault_detection} shows a consistent gradient trend across MR labels. Tests associated with \textit{Covered} relations detected 48 of 80 injected faults (60.0\%), compared with 38 of 80 (47.5\%) for \textit{Weak} relations, whereas none of the 40 faults associated with \textit{Untouched} relations were detected. The zero detection rate for \textit{Untouched} relations is expected: by definition, these relations correspond to behaviors that are not exercised by the original test suites. We include this group as a sanity-check baseline, rather than as evidence of relative oracle strength.

\begin{table}[h]
\centering
\caption{Detection of MR-relevant injected faults grouped by MR status.}
\label{tab:mr_fault_detection}
\small
\setlength{\tabcolsep}{4pt}
% \renewcommand{\arraystretch}{1.08}
\begin{tabular}{lrrr}
\toprule
\textbf{MR Status} & \textbf{\#Faults} & \textbf{\#Detected} & \textbf{Kill Rate} (\%) \\
\midrule
Covered   & 80 & 48 & 60.0 \\
Weak      & 80 & 38 & 47.5 \\
Untouched & 40 &  0 & 0.0 \\
\bottomrule
\end{tabular}
\end{table}

This pattern is consistent with the intended interpretation of MR labels: relations with explicit validation provide stronger protection against MR-relevant injected faults than relations that are only exercised, while unexercised relations provide no detection evidence in this study. A representative pair illustrates this distinction. A \textit{Covered} MR in the \textit{Infinite Scroll} component detected a fault replacing \textit{delay} with \textit{9999}, while a \textit{Weak} MR in the \textit{Space} component left all tests passing after changing \textit{wrap=false} to \textit{wrap=undefined}. Overall, these results provide supporting evidence that MR labels are consistent with differences in detecting MR-relevant behavioral faults in this targeted sample.


\begin{tcolorbox}[colback=gray!5, colframe=gray!30, title=\textbf{\small Answer to RQ4}, coltitle=black!90, left=2.5pt, top=3pt, bottom=3pt, right=2.5pt, boxsep=1.8pt]
Reported issue descriptions are substantially represented in the inferred MR space. Weak-oracle relations often identify exercised behaviors with limited oracle evidence, and MR labels show a consistent trend in the targeted injected-fault study. Overall, these results support the practical relevance of MR coverage as a behavioral validation signal.
\end{tcolorbox}

